\documentclass[10pt,a4paper]{amsart}
\usepackage{amsmath,amssymb,amsthm,graphicx}
% Retained preamble for the Stacks Project chapter snapshot in this directory.
% The chapter's own preamble.tex is a copy of the shared Stacks Project preamble; the block
% below is sliced MECHANICALLY, line for line, from it. Nothing was added, removed or
% redefined; the wrapper class, the xr-hyper/hyperref wiring and the externaldocument list
% are omitted because the assembler supplies the document class.
%
% lines taken from preamble.tex: 3, 5, 9, 11, 12, 156, 157, 158, ... (51 lines total)
%
% package coverage restored from the chapter's own preamble (2026-10-13): xy, xy-2cell

\usepackage[all]{xy}
\xyoption{2cell}
\UseAllTwocells

\newenvironment{reference}{\comment}{\endcomment}
\newenvironment{slogan}{\comment}{\endcomment}
\newenvironment{history}{\comment}{\endcomment}
\theoremstyle{plain}
\newtheorem{theorem}[subsection]{Theorem}
\newtheorem{proposition}[subsection]{Proposition}
\newtheorem{lemma}[subsection]{Lemma}
\theoremstyle{definition}
\newtheorem{definition}[subsection]{Definition}
\newtheorem{example}[subsection]{Example}
\newtheorem{exercise}[subsection]{Exercise}
\newtheorem{situation}[subsection]{Situation}
\theoremstyle{remark}
\newtheorem{remark}[subsection]{Remark}
\newtheorem{remarks}[subsection]{Remarks}
\numberwithin{equation}{subsection}
\def\lim{\mathop{\mathrm{lim}}\nolimits}
\def\colim{\mathop{\mathrm{colim}}\nolimits}
\def\Spec{\mathop{\mathrm{Spec}}}
\def\Hom{\mathop{\mathrm{Hom}}\nolimits}
\def\Ext{\mathop{\mathrm{Ext}}\nolimits}
\def\SheafHom{\mathop{\mathcal{H}\!\mathit{om}}\nolimits}
\def\SheafExt{\mathop{\mathcal{E}\!\mathit{xt}}\nolimits}
\def\Sch{\mathit{Sch}}
\def\Mor{\mathop{\mathrm{Mor}}\nolimits}
\def\Ob{\mathop{\mathrm{Ob}}\nolimits}
\def\Sh{\mathop{\mathit{Sh}}\nolimits}
\def\NL{\mathop{N\!L}\nolimits}
\def\CH{\mathop{\mathrm{CH}}\nolimits}
\def\proetale{{pro\text{-}\acute{e}tale}}
\def\etale{{\acute{e}tale}}
\def\QCoh{\mathit{QCoh}}
\def\Ker{\mathop{\mathrm{Ker}}}
\def\Im{\mathop{\mathrm{Im}}}
\def\Coker{\mathop{\mathrm{Coker}}}
\def\Coim{\mathop{\mathrm{Coim}}}
\DeclareMathSymbol{\boxtimes}{\mathbin}{AMSa}{"02}
\def\QCohstack{\mathcal{QC}\!\mathit{oh}}
\def\Cohstack{\mathcal{C}\!\mathit{oh}}
\def\Spacesstack{\mathcal{S}\!\mathit{paces}}
\def\Quotfunctor{\mathrm{Quot}}
\def\Hilbfunctor{\mathrm{Hilb}}
\def\Curvesstack{\mathcal{C}\!\mathit{urves}}
\def\Polarizedstack{\mathcal{P}\!\mathit{olarized}}
\def\Complexesstack{\mathcal{C}\!\mathit{omplexes}}
\def\Pic{\mathop{\mathrm{Pic}}\nolimits}
\def\Picardstack{\mathcal{P}\!\mathit{ic}}
\def\Picardfunctor{\mathrm{Pic}}
\def\Deformationcategory{\mathcal{D}\!\mathit{ef}}
\usepackage{cleveref}
\crefname{originallemma}{original lemma}{original lemmas}
\Crefname{originallemma}{Original lemma}{Original lemmas}
\makeatletter
\newlabel{algebraic-lemma-characterize-representable-by-algebraic-spaces}{{04Y5}{source}{}{sourcepointer-1}{}}
\newlabel{algebraic-lemma-characterize-representable-by-algebraic-spaces@cref}{{[originallemma][0][]04Y5}{[1][0][]source}{}{}{}}
\makeatother
\crefname{originallemma}{original lemma}{original lemmas}
\Crefname{originallemma}{Original lemma}{Original lemmas}
\makeatletter
\newlabel{categories-lemma-fully-faithful-diagonal-equivalence}{{04Z7}{source}{}{sourcepointer-2}{}}
\newlabel{categories-lemma-fully-faithful-diagonal-equivalence@cref}{{[originallemma][0][]04Z7}{[1][0][]source}{}{}{}}
\makeatother
\crefname{originallemma}{original lemma}{original lemmas}
\Crefname{originallemma}{Original lemma}{Original lemmas}
\makeatletter
\newlabel{lemma-check-property-covering}{{lemma-check-property-covering}{source}{}{sourcepointer-3}{}}
\newlabel{lemma-check-property-covering@cref}{{[originallemma][0][]lemma-check-property-covering}{[1][0][]source}{}{}{}}
\makeatother
\crefname{originallemma}{original lemma}{original lemmas}
\Crefname{originallemma}{Original lemma}{Original lemmas}
\makeatletter
\newlabel{lemma-check-representable-covering}{{lemma-check-representable-covering}{source}{}{sourcepointer-4}{}}
\newlabel{lemma-check-representable-covering@cref}{{[originallemma][0][]lemma-check-representable-covering}{[1][0][]source}{}{}{}}
\makeatother
\crefname{originallemma}{original lemma}{original lemmas}
\Crefname{originallemma}{Original lemma}{Original lemmas}
\makeatletter
\newlabel{section-properties-morphisms}{{section-properties-morphisms}{source}{}{sourcepointer-5}{}}
\newlabel{section-properties-morphisms@cref}{{[originallemma][0][]section-properties-morphisms}{[1][0][]source}{}{}{}}
\makeatother
\title{Monomorphism = fully faithful = diagonal equivalence = fppf-locally a monomorphism of spaces}
\author{The Stacks Project Authors}
\date{stacks-project:stacks-properties.tex:local-snapshot-20261012; GNU FDL 1.2 or later}
\begin{document}
\maketitle
\noindent\textit{Editorial note.} Statement span is lines 1673-1688 of sources/stacks-stacks-properties/pinned-original-source.tex; proof span is lines 1690-1730, the complete author proof (39 body lines).\par\medskip
\noindent\textit{Editorial note.} Navigation stubs: 5 cited-but-not-included labels are supplied as non-load-bearing pointers with their pinned tags where the snapshot tag index has them.\par\medskip
\noindent\textit{Editorial note.} extra\_packages is empty on purpose: the retained preamble already loads xy with options and a second declaration clashes in the assembler wrapper.\par\medskip
\noindent\textit{Editorial note.} Text is retained verbatim, including the author's own typos and placeholders.\par\medskip
\noindent\textit{Editorial note.} Queue index 8 of queue\_stacks-properties.json (pb=1698, stmt 1673-1688, proof 1690-1730).\par\medskip
\noindent\textit{Editorial note.} stacks-properties wave-1: fresh adjudication accept (2026-10-14); pool scanned while CONTROL.md is down (0 bytes) - §2 row deferred per Chief, claim reported to Chief 08:0x.\par\medskip
\noindent\textit{Allowed standard prerequisites (editorial).}\begin{itemize}
\item Explicit prerequisite boundary. The cited inputs are prerequisites; the argument in the retained span is not replaced by any cited result.
\end{itemize}\medskip
\ifdefined\hypertarget\hypertarget{sourcepointer-1}{}\fi
\noindent\textit{Source reference.} Original lemma 04Y5: sources/stacks-stacks-properties/pinned-original-source.tex, source lines 1730--1730. Cited by the retained statement or proof but not part of the collected spans; supplied only as a navigation pointer. The argument does not depend on its content.\par\medskip
\ifdefined\hypertarget\hypertarget{sourcepointer-2}{}\fi
\noindent\textit{Source reference.} Original lemma 04Z7: sources/stacks-stacks-properties/pinned-original-source.tex, source lines 1730--1730. Cited by the retained statement or proof but not part of the collected spans; supplied only as a navigation pointer. The argument does not depend on its content.\par\medskip
\ifdefined\hypertarget\hypertarget{sourcepointer-3}{}\fi
\noindent\textit{Source reference.} Original lemma lemma-check-property-covering: sources/stacks-stacks-properties/pinned-original-source.tex, source lines 683--693. Cited by the retained statement or proof but not part of the collected spans; supplied only as a navigation pointer. The argument does not depend on its content.\par\medskip
\ifdefined\hypertarget\hypertarget{sourcepointer-4}{}\fi
\noindent\textit{Source reference.} Original lemma lemma-check-representable-covering: sources/stacks-stacks-properties/pinned-original-source.tex, source lines 394--401. Cited by the retained statement or proof but not part of the collected spans; supplied only as a navigation pointer. The argument does not depend on its content.\par\medskip
\ifdefined\hypertarget\hypertarget{sourcepointer-5}{}\fi
\noindent\textit{Source reference.} Original lemma section-properties-morphisms: sources/stacks-stacks-properties/pinned-original-source.tex, source lines 378--401. Cited by the retained statement or proof but not part of the collected spans; supplied only as a navigation pointer. The argument does not depend on its content.\par\medskip
\begin{lemma}
\label{lemma-monomorphism}
Let $f : \mathcal{X} \to \mathcal{Y}$ be a morphism of algebraic stacks.
The following are equivalent:
\begin{enumerate}
\item $f$ is a monomorphism,
\item $f$ is fully faithful,
\item the diagonal
$\Delta_f : \mathcal{X} \to \mathcal{X} \times_\mathcal{Y} \mathcal{X}$
is an equivalence, and
\item there exists an algebraic space $W$ and a surjective, flat morphism
$W \to \mathcal{Y}$ which is locally of finite presentation such that
$V = \mathcal{X} \times_\mathcal{Y} W$ is an algebraic space, and the
morphism $V \to W$ is a monomorphism of algebraic spaces.
\end{enumerate}
\end{lemma}

\begin{proof}
The equivalence of (1) and (4) follows from the general discussion in
Section \ref{section-properties-morphisms}
and in particular
Lemmas \ref{lemma-check-representable-covering} and
\ref{lemma-check-property-covering}.

\medskip\noindent
The equivalence of (2) and (3) is
Categories, Lemma \ref{categories-lemma-fully-faithful-diagonal-equivalence}.

\medskip\noindent
Assume the equivalent conditions (2) and (3). Then $f$ is representable
by algebraic spaces according to
Algebraic Stacks,
Lemma \ref{algebraic-lemma-characterize-representable-by-algebraic-spaces}.
Moreover, the $2$-Yoneda lemma combined with the fully faithfulness
implies that for every scheme $T$ the functor
$$
\Mor(T, \mathcal{X})
\longrightarrow
\Mor(T, \mathcal{Y})
$$
is fully faithful. Hence given a morphism $y : T \to \mathcal{Y}$ there exists
up to unique $2$-isomorphism at most one morphism $x : T \to \mathcal{X}$
such that $y \cong f \circ x$. In particular, given a morphism of schemes
$h : T' \to T$ there exists at most one lift
$\tilde h : T' \to T \times_\mathcal{Y} \mathcal{X}$ of $h$.
Thus $T \times_\mathcal{Y} \mathcal{X} \to T$ is a monomorphism of
algebraic spaces, which proves that (1) holds.

\medskip\noindent
Finally, assume that (1) holds. Then for any scheme $T$ and morphism
$y : T \to \mathcal{Y}$ the fibre product $T \times_\mathcal{Y} \mathcal{X}$
is an algebraic space, and $T \times_\mathcal{Y} \mathcal{X} \to T$
is a monomorphism. Hence there exists up to unique isomorphism
exactly one pair $(x, \alpha)$ where $x : T \to \mathcal{X}$ is a morphism
and $\alpha : f \circ x \to y$ is a $2$-morphism. Applying
the $2$-Yoneda lemma this says exactly that $f$ is fully faithful, i.e.,
that (2) holds.
\end{proof}

\end{document}
